\documentclass[UTF8,10pt,fontset=none]{ctexart}
\usepackage[a4paper,margin=18mm]{geometry}
\usepackage{amsmath,amssymb,mathtools}
\usepackage{enumitem}
\setCJKmainfont{Songti SC}
\setmainfont{Times New Roman}
\setlength{\parindent}{0pt}
\setlist[enumerate]{leftmargin=2.3em,itemsep=.55em,topsep=.4em}
\pagestyle{empty}
\begin{document}
\begin{center}
  {\Large\bfseries 浙江大学 2024 年直博数学试题}
\end{center}
\vspace{1em}

\section*{一、数学分析}
\begin{enumerate}
\item 请利用有限覆盖定理证明闭区间上连续函数必一致连续。
\item 设可微函数列 $\{f_n(x)\}$ 在 $[0,1]$ 上处处收敛，且 $\{f'_n(x)\}$ 在 $[0,1]$ 上一致有界。证明：$\{f_n(x)\}$ 在 $[0,1]$ 上一致收敛。
\item 若 $f$ 在 $[0,+\infty)$ 上可导，且 $f'$ 在 $[0,+\infty)$ 上一致连续，且 $\lim_{x\to+\infty}f(x)$ 存在，证明：
\[
\lim_{x\to+\infty}f'(x)=0.
\]
\item 设函数 $f$ 在区间 $(0,1)$ 上连续。若对任意 $x\in(0,1)$，有
\[
\lim_{h\to0}\frac{f(x+h)-f(x-h)}{h}>0,
\]
则 $f$ 在 $(0,1)$ 上严格单调递增。
\item 设 $f$ 是以 $\pi$ 为周期的偶函数，且在 $[0,\pi]$ 上可积。证明：
\[
\int_0^{+\infty}f(x)\left(\frac{\sin x}{x}\right)^2\,dx
=\int_0^{\pi/2}f(x)\,dx.
\]
\end{enumerate}

\section*{二、高等代数}
\begin{enumerate}
\item 假如 $T$ 是实数域上 $n$ 维线性空间 $V$ 上的线性变换，满足 $T^2+I=0$，定义 $(a+bi)v=av+bTv$。证明 $V$ 关于上面的运算构成复数域上的线性空间，并求其在复数域上的维数。
\item 假如 $n$ 阶复矩阵 $A,B,C$ 满足 $AB-BA=\overline C C^{\mathrm T}$，其中 $\overline C$ 是 $C$ 的共轭矩阵，$C^{\mathrm T}$ 是 $C$ 的转置矩阵；两个复数 $a,b$ 满足 $f(a,b)=0$ 对所有满足 $f(A,B)=0$ 的多项式 $f(x,y)$ 都成立。证明存在非零向量 $v$ 满足 $Av=av$ 且 $Bv=bv$。
\item 假如 $f(x)$ 和 $g(x)$ 是两个互素的多项式，$A$ 是一个 $n$ 阶方阵，证明 $f(A)g(A)=0$ 的充分必要条件是 $f(A)$ 的秩与 $g(A)$ 的秩之和等于 $n$。
\item 假如 $V$ 是次数不大于 $2$ 的实系数多项式组成的线性空间，在 $V$ 中求多项式 $q(x)$，使得对 $V$ 中的所有 $p(x)$，等式
\[
\int_{-1}^{1}p(x)q(x)\,dx=\int_{-1}^{1}p(x)\cos(\pi x)\,dx
\]
都成立。
\item 已知 $n$ 阶行列式 $A=(a_{ij})$，$n\ge3$，证明 $\det A$ 是 $3$ 的整数倍当且仅当 $n$ 是奇数，其中
\[
a_{ij}=\begin{cases}
3i,&i=j,\\
3i+1,&i<j,\\
2-3j,&i>j.
\end{cases}
\]
\end{enumerate}
\end{document}
